\section{What a game specification must preserve}
\label{sec:semantics}

The standard nonlocal game is a cooperative test: a referee samples questions,
sends each player only their own question, receives answers, and evaluates a
score. Its meaning depends on the absence of communication during the round and
on what resources were distributed beforehand \cite{cleve2004}. Entanglement is
a resource that can generate Bell-nonlocal correlations; entanglement and Bell
nonlocality are not interchangeable properties. CHSH violation certifies
entanglement under its assumptions, but failure to violate a particular Bell
inequality does not certify separability.

\begin{definition}[Finite classical one-round verifier]
An arena is a finite list of parties $i\in\{1,\ldots,n\}$, each with a nonempty
question alphabet $X_i$ and answer alphabet $A_i$. Write
$X=\prod_iX_i$ and $A=\prod_iA_i$. A verifier is a distribution
$\mu:X\to[0,1]$ and a real payoff $V:A\times X\to\R$. A winning-probability
verifier has $0\leq V\leq1$; a deterministic win/lose predicate takes values in
$\{0,1\}$. A behavior is a normalized conditional distribution $p(a\mid x)$,
and its score is
\begin{equation}
 S_G(p)=\sum_{x\in X}\mu(x)\sum_{a\in A}V(a,x)p(a\mid x).
 \label{eq:score}
\end{equation}
The standard information rule gives player $i$ only $x_i$, with no interparty
communication during the round. Question sampling is independent of the
players' pre-existing resources.
\end{definition}

There is no invariant answer to ``how many fields define a game'': one may
encode all data in a single string, infer $n$ from an alphabet list, or combine
$\mu$ and $V$. A practical answer has five semantic components:
\begin{enumerate}
\item an arena with parties, interfaces, and their information/causal structure;
\item a verifier containing sampling, acceptance/payoff, and any trusted operations;
\item an admissible resource/strategy model $\mathcal R$;
\item an analysis question, such as a value, membership, separation, or extraction bound;
\item optional reference realizations and a target-equivalence specification.
\end{enumerate}
The first two specify the test. Adding the third defines the optimization
problem $\val_{\mathcal R}(G)=\sup_{p\in\mathcal R}S_G(p)$. Reference states
are not necessary to define $G$ or its value. They become essential when the
question is whether a score certifies a particular realization. The supremum
must not silently become a maximum.

Question-dependent answer alphabets can be represented using a common padded
alphabet only with explicit legality restrictions. Merely assigning zero payoff
to invalid answers is unsafe for signed payoffs: an invalid answer could become
an advantageous abort. The first implementation uses a fixed answer alphabet
for each party. A party with no external input has $|X_i|=1$; this is a syntax
convention, not permission to erase source-independence assumptions.

\subsection{An exact core representation, with a deliberate boundary}
Define $c_{a,x}=\mu(x)V(a,x)$. For a fixed arena,
$S_G(p)=\langle c,p\rangle$ is an exact intermediate representation (IR) for
linear score evaluation. Every finite classical one-round verifier has this
representation, and a bounded table can be serialized exactly when its entries
are rational. It does not follow that every physically meaningful experiment
is such a verifier, or that arbitrary real parameters have finite exact
encodings.

The implementation should retain both $\mu$ and $V$ alongside $c$. Different
sampling schemes can have the same expected linear score yet different
estimator ranges, variances, and experimental costs. Constraints on the
resource set also refer to the arena rather than just to coefficients. In
particular, a question tuple with zero sampling probability still belongs to
the arena; a non-signalling extension may impose constraints involving it.
Dropping those questions is a separate transformation requiring justification.

There are several inequivalent notions of ``the same game.'' Exact transcript
equivalence preserves the distribution and score of each run. Score
equivalence preserves $S_G(p)$ for every admissible $p$. Value equivalence
preserves only an optimum for one specified resource model. A relaxation
preserves neither behavior nor optimum, but supplies an inequality. A
compiler pass should state which of these contracts it provides.

A coefficient table costs $|X||A|$ scalar entries. A compact parity predicate
can be exponentially smaller. The universal finite table is therefore a
\emph{semantic fallback}, not a zero-cost implementation. Structured
expressions---parity tests, consistency constraints, graph predicates, and
local wirings---should survive until a backend actually needs expansion.
